\documentclass[11pt]{article}
\usepackage[T1]{fontenc}
\usepackage[a4paper,margin=1in]{geometry}
\usepackage{amsmath,amssymb,amsthm}
\usepackage[hidelinks]{hyperref}
\usepackage{microtype}

\newtheorem*{question}{Source Question 5.2}
\newtheorem*{answer}{Later answer}

\setlength{\parindent}{0pt}
\setlength{\parskip}{0.6em}

\title{Literature Status: Removing the HCR Hypothesis\\
for Cyclic Reproducing Kernels in $\mathcal H(b)$}
\author{Run \texttt{fa\_banach\_001} --- lane 7}
\date{11 August 2026}

\begin{document}
\maketitle

\begin{center}
\fbox{\parbox{0.9\textwidth}{
\textbf{Status: literature already answered.}
Question~5.2 of arXiv:2210.15735 asks whether the HCR/corona-pair
hypothesis in its Proposition~5.7 can be removed.  Theorem~4.2,
Corollary~5.1, and Example~5.2 of the separate later paper
arXiv:2403.04335 answer both parts affirmatively.  The later authors
explicitly state that the earlier corona-pair assumption can be omitted.
}}
\end{center}

\section{Original question}

Let $b$ be a non-extreme point of the closed unit ball of $H^\infty$, let
$a$ be its Pythagorean mate, and let $S_b$ denote the forward shift on the
de Branges--Rovnyak space $\mathcal H(b)$.  Fricain and Grivaux prove in
Proposition~5.7 of \cite{source}, under their condition (HCR), that every
reproducing kernel
\[
 k^b_\lambda(z)=\frac{1-\overline{b(\lambda)}b(z)}
                         {1-\overline\lambda z}
\]
is cyclic for $S_b$, and that if $b$ is outer then $bk_\lambda$ is cyclic
for every $\lambda\in\mathbb D$ (hence $b$ itself is cyclic).

On arXiv PDF page~20, the paper ends with:

\begin{question}
Does Proposition~5.7 hold without the assumption that $(a,b)$ satisfies
(HCR)?
\end{question}

\section{Separate later answer}

Fricain and Lebreton prove the following sufficient condition in
Theorem~4.2 of \cite{answerpaper}: if $f\in\mathcal H(b)$ is outer and
$b/f\in L^\infty(\mathbb T)$, then $f$ is cyclic for $S_b$.  Their
Corollary~5.1 gives two immediate consequences:
\begin{enumerate}
\item if $\inf_{\mathbb D}|f|>0$, then $f$ is cyclic;
\item $b$ is cyclic if and only if $b$ is outer.
\end{enumerate}

Example~5.2 on arXiv PDF page~16 applies these statements to the exact two
parts of Proposition~5.7.  For each $\lambda\in\mathbb D$,
\[
 |k^b_\lambda(z)|
 =\left|\frac{1-\overline{b(\lambda)}b(z)}{1-\overline\lambda z}\right|
 \geq \frac{1-|b(\lambda)|}{2}>0,
\]
so Corollary~5.1(i) makes $k^b_\lambda$ cyclic.  If $b$ is outer, then
$bk_\lambda$ is outer and
\[
 \frac{b}{bk_\lambda}=\frac1{k_\lambda}=1-\overline\lambda z
 \in L^\infty(\mathbb T),
\]
so Theorem~4.2 makes $bk_\lambda$ cyclic.  Corollary~5.1(ii) also gives
the stated cyclicity of $b$ itself.

The authors then explicitly note that these were the two results previously
obtained in Proposition~5.7 under the corona-pair assumption, and conclude
that the assumption can be omitted.  Thus Question~5.2 has a complete
affirmative answer.

\section{Scope}

The identification is exact: it covers every non-extreme $b$, every
$\lambda\in\mathbb D$, and both conclusions of the source proposition.
No new mathematical result is claimed in this packet.

\begin{thebibliography}{9}
\bibitem{source}
E. Fricain and S. Grivaux,
\emph{Cyclicity in de Branges--Rovnyak spaces},
arXiv:2210.15735.

\bibitem{answerpaper}
E. Fricain and R. Lebreton,
\emph{Cyclicity of the shift operator and a related completeness problem
in de Branges--Rovnyak spaces},
arXiv:2403.04335.
\end{thebibliography}

\end{document}
